\section{Conclusions and Future Work}
\label{sec:conclusions}
\vspace{1pt}

We introduced TarMAC, an architecture for multi-agent reinforcement learning that
allows targeted continuous communication between agents via a sender-receiver soft attention
mechanism and multiple rounds of collaborative reasoning. Evaluation on four
diverse environments shows that our model is able to learn intuitive communication
attention behavior and improves performance, even in non-cooperative settings,
with task reward as sole supervision. While TarMAC uses continuous vectors as messages,
it is possible to force these to be discrete, either during training itself
(as in~\citet{foerster_nips16}) or by adding a decoder after learning to ground
these messages into symbols.

In future, we aim to exhaustively benchmark TarMAC on more challenging
$3$D navigation tasks because we believe this is where decentralized targeted
communication is most crucial, as it allows scaling to a large number of agents
with high-dimensional observation spaces. In particular, we are interested
in investigating combinations of TarMAC with recent advances in spatial memory,
planning networks, \etc.

